%% kapa: exam=1 round=27 year=2016 period=2 form=A subject=appraisal_law subject_ko=감정평가관계법규 seq=23 number=23
%% source: Q-Net 감정평가사 기출문제 (cst00309) artlSeq=5169402 file=extracted/1008/5169402/제27회 감정평가사 1차 문제/제27회 감정평가사 2교시문제(A형).pdf
%% conversion: kordoc 4.24.51 -> tools/gen1.py (auto); review=proofread-sonnet; flags=none
\begin{question}{23}
국유재산법령상 부동산인 행정재산의 사용료에 관한 설명으로 옳지 않은 것은?
\begin{choices}
  \item 경쟁입찰로 사용허가를 하는 경우 첫 해의 사용료는 최고입찰가로 결정한다.
  \item 사용료가 100만원을 초과하는 경우에 연 6회 이내에서 나누어 내게 할 수 있다.
  \item 사용료를 나누어 내게 할 때 연간 사용료가 1천만원 이상인 경우에는 그 허가를 받은 자에게 연간 사용료의 100분의 50에 해당하는 금액의 범위에서 보증금을 예치하게 하거나 이행보증조치를 하도록 하여야 한다.
  \item 중앙관서의 장은 행정재산을 공용으로 사용하려는 지방자치단체에 사용허가하는 경우 사용료를 면제하여야 한다.
  \item 보존용재산을 사용허가하는 경우에 재산의 유지ㆍ보존을 위하여 관리비가 특히 필요할 때에는 사용료에서 그 관리비 상당액을 뺀 나머지 금액을 징수할 수 있다.
\end{choices}
\end{question}
\answer{4}
